




\subsection{Модуль lab3}

\subsubsection{FlowNetwork: сеть потоков}

\textbf{Вход:} флаг ориентированности (конструктор, по умолчанию
\texttt{true}).

\textbf{Выход:} \texttt{unique\_ptr<FlowNetwork>}.

\textbf{Описание:} \texttt{FlowNetwork} наследует
\texttt{GraphBase<FlowVertex, FlowEdge>}. Структура \texttt{FlowEdge}
хранит \texttt{from}, \texttt{to}, \texttt{capacity}, \texttt{cost}
и \texttt{flow}. При добавлении ребра через \texttt{addEdge}
автоматически создаётся обратное ребро с нулевой пропускной
способностью и отрицательной стоимостью~--- это стандартное
требование алгоритмов на остаточных сетях.
Остаточная пропускная способность вычисляется как
$\mathit{res}(u,v) = \mathit{capacity}(u,v) - \mathit{flow}(u,v)$.
Метод \texttt{addFlow} обновляет поток симметрично: увеличивает
прямое ребро и уменьшает обратное.

\begin{lstlisting}[style=cppstyle, caption={Добавление ребра и обратного ребра (FlowNetwork.cc)}]
bool FlowNetwork::addEdge(int from, int to,
                          double capacity, double cost)
{
    FlowEdge forward_edge(from, to, capacity, cost, 0.0);
    FlowEdge backward_edge(to, from, 0.0, -cost, 0.0);

    bool added = GraphBase::addEdge(from, to, forward_edge);
    if (added) {
        GraphBase::addEdge(to, from, backward_edge);
        m_vertices_[to]->addNeighbor(from);
    }
    return added;
}

void FlowNetwork::addFlow(int from, int to, double flow) {
    auto* edge_forward  = getEdgeMutable(from, to);
    auto* edge_backward = getEdgeMutable(to, from);
    if (edge_forward)  edge_forward->flow  += flow;
    if (edge_backward) edge_backward->flow -= flow;
}
\end{lstlisting}

\subsubsection{MaxFlow: алгоритм Форда--Фалкерсона}

\textbf{Вход:} \texttt{FlowNetwork\&} (конструктор), \texttt{int source},
\texttt{int sink}, \\ \texttt{bool enableLogging}.

\textbf{Выход:} \texttt{double}~--- значение максимального потока.
Побочный эффект~--- вектор \texttt{m\_snapshots\_} для анимации.

\textbf{Описание:} Реализован метод Форда--Фалкерсона с BFS
для поиска увеличивающего пути.
На каждой итерации BFS ищет путь в остаточной сети~---
ребро рассматривается только если
$\mathit{res}(u,v) > 0$. Пропускная способность пути~---
минимум остаточных пропускных способностей по всем рёбрам~---
вычисляется через \texttt{PathUtils::getMinPathValue}.
Поток увеличивается через \texttt{PathUtils::forEachEdgeInPath}.
При включённом логировании после каждой итерации фиксируется
снимок состояния сети \texttt{FlowSnapshot}~--- он используется
для построения GIF-анимации роста потока.
Сложность: $O(V \cdot E^2)$.

\begin{lstlisting}[style=cppstyle, caption={Алгоритм Форда--Фалкерсона (MaxFlow.cc)}]
double MaxFlow::fordFulkerson(int source, int sink,
                              bool enableLogging)
{
    double max_flow = 0.0;
    std::unordered_map<int, int> parent;
    int step = 0;

    resetFlows(m_network_);
    m_snapshots_.clear();
    if (enableLogging) captureSnapshot(step++, max_flow);

    while (bfs(source, sink, parent)) {
        auto calc_residual = [this](int u, int v) {
            return m_network_.getResidualCapacity(u, v);
        };
        double path_flow = PathUtils<double>::getMinPathValue(
            source, sink, parent, calc_residual);

        PathUtils<double>::forEachEdgeInPath(
            source, sink, parent,
            [this, path_flow](int u, int v) {
                m_network_.addFlow(u, v, path_flow);
            });

        max_flow += path_flow;
        if (enableLogging) captureSnapshot(step++, max_flow);
    }
    return max_flow;
}
\end{lstlisting}

\subsubsection{MinCostFlow: поток минимальной стоимости}

\textbf{Вход:} \texttt{FlowNetwork\&} (конструктор), \texttt{int source},
\texttt{int sink}, \texttt{double targetFlow}.

\textbf{Выход:} \texttt{Result}~--- структура с полями \texttt{flow},
\texttt{cost}, \texttt{success} и \texttt{path} последнего найденного пути.

\textbf{Описание:} Алгоритм последовательно ищет кратчайший путь
в остаточной сети по стоимости через алгоритм Беллмана--Форда
и направляет по нему максимально возможный поток, не превышающий
$\mathit{targetFlow} - \mathit{currentFlow}$. Беллман--Форд
корректно обрабатывает отрицательные стоимости обратных рёбер.
По условию задания целевой поток берётся как $\lfloor 2/3 \cdot
\mathit{maxFlow} \rfloor$. Сложность одной итерации: $O(V \cdot E)$.

\begin{lstlisting}[style=cppstyle, caption={Поток минимальной стоимости (MinCostFlow.cc)}]
MinCostFlow::Result MinCostFlow::findMinCostFlow(
    int source, int sink, int targetFlow)
{
    Result result;
    resetFlows(m_network_);
    double current_flow = 0.0, total_cost = 0.0;
    std::unordered_map<int, double> dist;
    std::unordered_map<int, int>    parent;

    while (current_flow < targetFlow
           && bellmanFord(source, sink, dist, parent))
    {
        auto calc_residual = [this](int u, int v) {
            return m_network_.getResidualCapacity(u, v);
        };
        double path_flow = std::min(
            targetFlow - current_flow,
            PathUtils<double>::getMinPathValue(
                source, sink, parent, calc_residual));

        result.path = PathUtils<double>::reconstructPath(
            source, sink, parent);

        PathUtils<double>::forEachEdgeInPath(
            source, sink, parent,
            [this, &total_cost, path_flow](int u, int v) {
                m_network_.addFlow(u, v, path_flow);
                total_cost += path_flow
                            * m_network_.getCost(u, v);
            });

        current_flow += path_flow;
    }
    result.flow    = current_flow;
    result.cost    = total_cost;
    result.success = (current_flow == targetFlow);
    return result;
}
\end{lstlisting}

\subsubsection{PathUtils: вспомогательные операции на путях}

\texttt{PathUtils<T>}~--- шаблонный статический класс для работы
с путями, заданными массивом предшественников
\texttt{unordered\_map<int,int> parent}. Используется как в
\texttt{MaxFlow}, так и в \texttt{MinCostFlow}, а также в
\texttt{Dijkstra} из лабораторной~2.

\texttt{getMinPathValue} вычисляет минимальное значение геттера
по всем рёбрам пути~--- используется для нахождения
пропускной способности пути. \texttt{forEachEdgeInPath} применяет
произвольный функтор к каждому ребру пути~--- использует
